\documentclass[10pt]{article}
\usepackage[margin=20mm]{geometry}
\usepackage{amsmath,amssymb,mathtools,fontspec,unicode-math,xurl,hyperref}
\IfFontExistsTF{DejaVu Serif}{\setmainfont{DejaVu Serif}}{\setmainfont{Latin Modern Roman}}
\setmathfont{latinmodern-math.otf}
\setlength{\parindent}{0pt}\setlength{\parskip}{6pt}\sloppy
\title{Initial Scientific Review of
Convergence Properties of Dynamic Mode
Decomposition for Analytic Interval Maps}
\author{Bounded scientific review: M25-0425}
\date{}
\begin{document}
\maketitle
Abstract: The acquired 24-page manuscript contains a false boundedness assertion in the analytic transfer-operator factorization. An explicit expanding full-branch interval map satisfies even the stronger expansion condition used in the main theorem, but its transfer operator does not map the asserted \(H^\infty\) spaces into one another. This is a Major problem in the derivation of the quantitative convergence estimate. It does not establish a counterexample to the final spectral-convergence conclusion. An isolated typographical error in the Hardy coefficient norm is retained as an uncounted editorial clarification. The finite-sample conditioning argument, polynomial example, and unperturbed Fourier example were also checked. Additional restrictions needed to interpret certain proof steps and numerical examples are stated without promoting them to unproved counterexamples. Scientific adjudication remains pending.\par
\section*{Identity and scope}
The source is E. Akindji, J. Slipantschuk, O. F. Bandtlow, and W. Just, Convergence properties of dynamic mode decomposition for analytic interval maps, DOI https://doi.org/10.1002/cpa.70011. The acquired PDF is arXiv:2404.08512v1, deposited 12 April 2024, with a manuscript date of 10 April 2024 [source]. It has 24 pages and SHA256 `932f2146b632f41c9734c842116236e242665056b9e8927a6c27957e67e4f851`. Every extracted source page, including the numerical discussion and references, was read. Numbered locators refer to this exact version; physical and printed page numbers agree.\par
The annual task’s DOI identity is retained. No equivalence between this preprint and a later journal version is asserted. The retained source intake does not supply an author-code implementation or a supplementary proof document; the bounded recovery search found no substantive earlier report for this paper. This is therefore an initial review of the acquired source. Missing code is not treated as a mathematical defect, and the published numerical figures were not independently regenerated.\par
\section*{Framework and checks}
Let \(D_s=\{z\in\mathbb C:|z|<s\}\). The standing assumptions on p. 5 specify finitely many analytic inverse branches \(\phi_\ell\in H^\infty(D_R)\), with \(\phi_\ell(D_R)\subset D_r\) for \(1<r<R\). The transfer operator is \[(\mathcal Lf)(z)=\sum_{\ell=1}^d
 \sigma_\ell\phi_\ell'(z)f(\phi_\ell(z)),\qquad
 \sigma_\ell=\operatorname{sgn}\phi_\ell'(0).\] The signs encode the real branch orientations, and the change of variables gives the bilinear duality in (5). The auxiliary Hardy space has the radial boundary norm and normalized monomials \(e_n(z)=(z/\rho)^n\). The main theorem imposes \[\gamma=(1+\sqrt2)^2,\qquad \gamma r<R,\qquad
 \rho_*=\sqrt{rR/\gamma}.\] Its matrix estimate, for sufficiently many equidistant nodes, is \[\left\lVert V_N(H_N^{(M)})^\dagger G_N^{(M)}V_N^{-1}-L\right\rVert
 \le C\left\{\frac{(\gamma rR)^{N/2}N^2}{M}
                 +\left(\frac{\gamma r}{R}\right)^{N/2}\right\}.\] Here \(L\) is the matrix representation of \(\mathcal L\) on the chosen Hardy space; \(V_N\) changes between monomials and the normalized basis. The finite-section convention matters because \(V_N^{-1}\) is understood on that section.\par
The Riemann-sum argument and Schur estimate in Lemmas 2.5–2.6 and Proposition 2.7 give the stated order \(N^2/M\). The Gram matrix is positive definite when \(M\ge N\): a nonzero polynomial of degree at most \(N-1\) cannot vanish at all \(M\) distinct nodes. The perturbation argument of Lemma 2.12 then combines the inverse Hilbert-matrix growth \(C\gamma^N\) with the quadrature error and a Neumann-series bound. No failure was established in this conditioning step. The cited sharp Hilbert-matrix asymptotics were used as an external lemma, not independently re-proved or numerically certified.\par
\section*{Supported problems}
Problem 1 (Major: unbounded derivative weights in Lemma 2.3).  Lemma 2.3, p. 7, asserts that \(\widetilde{\mathcal L}:H^\infty(D_r)\to H^\infty(D_R)\) is bounded. Its proof bounds the operator norm by \[B=\sup_{z\in D_R}\sum_{\ell=1}^d|\phi_\ell'(z)|.\] The standing assumptions do not make \(B\) finite. Boundedness of a holomorphic function on an open disk does not imply boundedness of its derivative on that same disk.\par
Here is an admissible example, including the real full-branch condition. Take \(R=10\), \(r=3/2\), and partition \([-1,1]\) into 101 equal intervals. For \(\ell=1,\ldots,101\), except \(\ell=51\), put \[c_\ell=\frac{2\ell-102}{101},\qquad
 \phi_\ell(z)=c_\ell+\frac{z}{101}.\] For the central interval define \[s_\pm=\sqrt{1\pm 1/10},\qquad
 k=\frac{2}{101(s_+-s_-)},\qquad
 \phi_*(z)=-\frac1{101}
       +k\left(s_+-\sqrt{1-z/10}\right).
 \tag{3.1}\] The principal square root is analytic on \(D_{10}\), since \(\operatorname{Re}(1-z/10)>0\). Each linear branch maps \([-1,1]\) increasingly onto its prescribed interval. The central branch has \(\phi_*(-1)=-1/101\), \(\phi_*(1)=1/101\), and \[\phi_*'(z)=\frac{k}{20\sqrt{1-z/10}}.
 \tag{3.2}\] It too is strictly increasing on the real interval. Rationalizing \(s_+-s_-\) gives \(k=10(s_++s_-)/101<20/101\). Hence \(0<\phi_*'(x)<10/909<1\) for \(-1\le x\le1\); all the inverse branches are contractions there. Their inverses define an expanding analytic full-branch map, with the usual one-sided endpoint convention.\par
For \(z\in D_{10}\), the noncentral branches satisfy \[|\phi_\ell(z)|<\frac{110}{101}<\frac32.\] For the central branch, using \(s_+<11/10\) and \(\sqrt2<3/2\), \[|\phi_*(z)|
 <\frac1{101}+\frac{20}{101}\left(\frac{11}{10}+\frac32\right)
 =\frac{53}{101}<\frac32.\] Thus every branch belongs to \(H^\infty(D_{10})\) and maps that disk into \(D_{3/2}\). The stronger main-theorem condition is also satisfied: \(\gamma r=9/2+3\sqrt2<9<10=R\).\par
All real orientation signs are positive. Applying the claimed operator to the admissible function \(f\equiv1\) gives \[(\widetilde{\mathcal L}1)(z)=\frac{100}{101}
                  +\frac{k}{20\sqrt{1-z/10}}.\] This tends to \(+\infty\) as real \(z\uparrow10\). Consequently \(\widetilde{\mathcal L}1\notin H^\infty(D_{10})\). There is no cancellation that could rescue the claim. This is a direct counterexample to Lemma 2.3, not merely a possibly infinite upper bound.\par
\paragraph*{Dependency and scope.}\mbox{}\par
Proposition 2.4, pp. 7–8, uses precisely this bounded map in \(\mathcal L=\mathcal J_2\widetilde{\mathcal L}\mathcal J_1\); its constant is \(\rho B/\sqrt{\rho^2-r^2}\). The finite-rank approximation bounds feed Proposition 2.11 and then Theorem 2.14 and Corollaries 2.15–2.17. The proof therefore lacks a finite constant at an essential step. This substantial dependency makes the problem Major even though there are natural repairs.\par
The example does not show that \(\mathcal L\) is unbounded on every smaller Hardy space, or that the final spectral conclusion fails. Indeed, choosing \(R'\) with \(\gamma r<R'<R\) bounds all branch derivatives on \(D_{R'}\) by Cauchy’s estimate and recovers the displayed factorization with \(R'\). The corresponding proof has new constants, the radius \(\sqrt{rR'/\gamma}\), and a rate involving \(R'\). Alternatively, adding \(\phi_\ell'\in H^\infty(D_R)\), or extension to a neighborhood of the closed disk, supplies the missing hypothesis. Shrinking the disk alone does not verify the quantitative statement with the original fixed \(R\). A revised proof or explicit strengthening is needed at that level.\par
\subsection*{Editorial clarification: the Hardy coefficient exponent}
After (6), the source prints \[\left\lVert f\right\rVert_{H^2(D_R)}^2=\sum_{n=0}^\infty |f_n|^2 R^{-2n}.\] For the Taylor coefficients defined immediately above, the radial boundary definition instead gives \[\left\lVert f\right\rVert_{H^2(D_R)}^2=\sum_{n=0}^\infty |f_n|^2 R^{2n}.\] Already \(f(z)=z\) and \(R>1\) distinguish the two values, \(R^2\) and \(R^{-2}\). The accompanying summability characterization needs the same correction.\par
This is local: the proof of Lemma 2.1 on the same page uses the correct positive exponent, as do Lemma 2.2 and the normalized basis later on. No separate failure of the main argument was derived from this typographical error. Replacing both negative exponents on p. 6 is the appropriate repair. This isolated typographical correction is not counted as a substantive Problem.\par
\section*{Additional proof and example boundaries}
\paragraph*{The radius in Proposition 2.13.}\mbox{}\par
The statement on p. 15 allows every \(\rho\in(r,R)\). In its proof the uniform bound on \(V_NH_N^\dagger G_NV_N^{-1}\) is obtained by quoting the optimized error \((\gamma r/R)^{N/2}\). Proposition 2.11 gives that optimized rate at \(\rho=\rho_*\); its general estimate includes \((\gamma\rho/R)^N\). The printed proof consequently does not justify the asserted uniform bound for an arbitrary \(\rho\) near \(R\). Restricting to \(\rho_*\), as used by Theorem 2.14, resolves this particular range issue. More generally \(r<\rho<R/\gamma\) is sufficient for the displayed boundedness argument, once Problem 1 is repaired. No counterexample to the broader proposition was established, so this observation is retained as a proof-range clarification.\par
The same proof has local transcription slips: the difference-of-inverses identity requires \(H_N-H_N^{(M)}\) in the correction term, and the later similarity expression needs \(V_N^{-1}\) on the right. Correcting the sign does not change the ensuing triangle estimate. The equality between a norm and the sum of two norms there should be an inequality. These observations do not constitute an additional independent Major finding.\par
\paragraph*{Fixed spectral modes.}\mbox{}\par
The standard perturbation argument behind Lemma 2.16 applies to a fixed isolated nonzero eigenvalue of the compact limit, with its finite algebraic multiplicity. If its largest Jordan block has size \(m\), a norm error \(O(e^{-\alpha N})\) can give an eigenvalue error of order \(O(e^{-\alpha N/m})\), still exponential for that fixed mode. This does not supply a uniform rate for the entire accumulating spectrum, nor does it by itself settle behavior at the non-isolated point zero. The Koopman identification through (5) and bounded functionals in (18) was checked; no counterexample to the paper’s fixed nonzero-mode conclusion was obtained.\par
\paragraph*{Polynomial and Fourier checks.}\mbox{}\par
For the skewed doubling map let \(p=(1+a)/2\), \(q=(1-a)/2\). Its inverse branches are \(py-q\) and \(qy+p\), so \[\mathcal Ly^n=p(py-q)^n+q(qy+p)^n.\] The polynomial subspaces are invariant, and the triangular diagonal is \(p^{n+1}+q^{n+1}\). Thus the infinite-node monomial EDMD matrices have exactly these finite-section eigenvalues. This confirms the mechanism explained in Section 3.1; it is not a finite-precision stability guarantee.\par
For centered Fourier indices \(k,\ell=-K,\ldots,K\), integration over the two real branches gives, with \(\operatorname{sinc}(t)=\sin(\pi t)/(\pi t)\), \[G_{k\ell}
 =p e^{i\pi\ell q}\operatorname{sinc}(k-p\ell)
  +q e^{-i\pi\ell p}\operatorname{sinc}(k-q\ell).\] At \(a=0\), this reduces to \(G_{k\ell}=(-1)^k\mathbf1_{\{\ell=2k\}}\). The constant mode has eigenvalue one; the remaining finite matrix is nilpotent. For \(K\ge1\) its largest nilpotent Jordan block has size \(\lfloor\log_2K\rfloor+1\). This supports the stated sensitivity mechanism near the unskewed case. The subsequent spectral-limit discussion is heuristic and was not converted into a proved universal convergence theorem here.\par
The two-branch examples do not instantiate the strong expansion hypothesis of Theorem 2.14. For the Blaschke inverse branches displayed on p. 18, \(\phi_\ell'(0)=1/2\). If \(\phi_\ell(D_R)\subset D_r\), Cauchy’s coefficient estimate implies \(r/R\ge1/2\), incompatible with \(r/R<1/\gamma=3-2\sqrt2\). The skewed linear example has the same obstruction because its largest inverse slope is at least \(1/2\). Their observed convergence is therefore evidence about a wider regime or special structure, not an application of that sufficient theorem. This does not refute the observed convergence. The nonlinear exact spectrum quoted from prior work was not independently re-proved, and the source’s plots were not reproduced.\par
\section*{Conclusion and evidence limits}
The review supports one Major problem and no Minor problems in the acquired preprint. The Hardy coefficient typo is retained as an uncounted editorial clarification. The decisive Major evidence is the fully specified full-branch counterexample to Lemma 2.3 and its explicit dependence in the main proof. The weaker conclusion that suitable smaller disks yield bounded derivatives was also checked, so the report does not overstate the consequence as a disproof of all intended convergence results.\par
The entire acquired text was read; the central source pages were additionally inspected as images. Native bounded helpers checked the analytic factorization and the example calculations. Their work is attributed in the retained evidence, separately from the primary review. Independent elementary checks are distinguished from author-code reproduction and from unverified external spectral results. All mathematical judgments here concern the exact acquired version. They do not certify every uncriticized statement, establish a correction in a later version, or confer independent acceptance.\par
Recorded dimension | State\par
Supported problems | 1 Major; 0 Minor\par
Scientific decision | Pending\par
Human reviews | None recorded\par
Independent acceptance | Not completed\par
1 E. Akindji, J. Slipantschuk, O. F. Bandtlow, and W. Just, Convergence properties of dynamic mode decomposition for analytic interval maps, arXiv:2404.08512v1, 12 April 2024, 24 pp. https://arxiv.org/abs/2404.08512v1. Task DOI: https://doi.org/10.1002/cpa.70011.\par
\end{document}
